(1) 线性化连续性与动量方程分别为 \(\rho'_t+\rho\nabla\cdot\boldsymbol u=0\)、\(\rho\boldsymbol u_t=-\nabla p'\)，状态关系 \(p'=c^2\rho'\)。设 \(\boldsymbol u=\nabla\phi\)、\(p'=-\rho\phi_t\)，消元得波动方程 \(\phi_{tt}=c^2\nabla^2\phi\)。取 \(e^{-i\omega t}\) 时间因子后，
\[
\widehat\phi''+\frac2r\widehat\phi'+k^2\widehat\phi=0.
\]
一般解为 \((Ae^{ikr}+Be^{-ikr})/r\)，无入射条件要求 \(B=0\)。因此
\[
\widehat u_r(r)=A\frac{(ikr-1)e^{ikr}}{r^2},\qquad
\widehat p(r)=i\rho\omega A\frac{e^{ikr}}r.
\]
在源面强制 \(\widehat u_r(a)=V_0\)，得到
\[
\boxed{A=\frac{V_0a^2e^{-ika}}{ika-1}}.
\]
代回即为任意半径的唯一出射解。必须保留 \(ika-1\) 这一近场导数因子，直接用远场 \(u=p/(\rho c)\) 代替源面边界会丢失球面储能及声阻抗的虚部。

(2) 令 \(x=ka\)，球面单位速度的复声压比为
\[
\boxed{Z_{\rm rad}=\frac{i\rho\omega a}{ix-1}
=\rho c\frac{x^2-ix}{1+x^2}}.
\]
因此 \(\operatorname{Re}Z=\rho c x^2/(1+x^2)\ge0\)，\(\operatorname{Im}Z=-\rho cx/(1+x^2)\)。前者描述每周期净输出的声能，后者代表近场中储能返回源的反应分量而非负耗散。低频
\(Z=\rho c[-ix+x^2+O(x^3)]\)，其中净辐射实部比反应性虚部小一个 \(x\) 阶；高频 \(Z\to\rho c\)，趋于远离源的平面波声阻抗。虚部符号依赖所指定的 \(e^{-i\omega t}\) 约定，不能把符号变化误当作吸收或增益。

(3) 复声强定义 \(\mathscr I_c=\widehat p\widehat u_r^*/2\)，由第 (1) 问逐项计算
\[
\mathscr I_c(r)=\frac{\rho\omega|A|^2}{2r^3}(kr-i).
\]
故时间平均真实辐射强度为
\(\overline I_r=\operatorname{Re}\mathscr I_c=\rho\omega k|A|^2/(2r^2)\)，总功率是
\[
\boxed{\overline P=4\pi r^2\overline I_r
=2\pi\rho c a^2|V_0|^2\frac{(ka)^2}{1+(ka)^2}},
\]
与观测半径无关。声源表面所做功率 \(4\pi a^2 \operatorname{Re}Z_{\rm rad}|V_0|^2/2\) 与此完全相等，直接验证出射能量守恒。在 \(ka\ll1\)、\(Q_0=4\pi a^2V_0\) 固定时，
\[
\boxed{\overline P\simeq\frac{\rho\omega^2|Q_0|^2}{8\pi c}}.
\]
而如果固定的是 \(V_0\)，同一式写作 \(\overline P\propto\omega^2a^4|V_0|^2\)，必须区分固定边界速度和固定体积速度的不同极限。

(4) 当 \(kr\gg1\)，有 \(\widehat u_r\sim ikA e^{ikr}/r\) 与 \(\widehat p\sim i\rho\omega A e^{ikr}/r=\rho c\,\widehat u_r\)，两者同相且按 \(r^{-1}\) 衰减；当 \(kr\ll1\)，速度主项为 \(-A/r^2\)，声压仍为 \(i\rho\omega A/r\)，两者近乎正交。复声强虚部为
\[
\boxed{\operatorname{Im}\mathscr I_c=-\frac{\rho\omega|A|^2}{2r^3}},
\]
比真实声强更快衰减，不构成一个与观测半径独立的净传播功率。固定 \(Q_0\)，使 \(a\to0\) 时 \(A\to-Q_0/(4\pi)\)；于是外场势收敛到 \(\widehat\phi=-Q_0e^{ikr}/(4\pi r)\)，得到点单极子辐射解。但呼吸球表面速度此时 \(V_0=Q_0/(4\pi a^2)\to\infty\)，最终破坏低马赫数、小变形的线性声学假设，故点单极是外场数学极限而非真实无限小球的物理操作说明。